
\subsection{Corrections des Exercices}

\begin{correctionbox}[Correction Exercice 1 : Concepts (Moments)]
1.  Le 1er moment non centré $E[X^1]$ est l'espérance $\mu$.
2.  Le 2ème moment centré $E[(X-\mu)^2]$ est la variance $\sigma^2$.
\end{correctionbox}

\begin{correctionbox}[Correction Exercice 2 : Interprétation (Skewness)]
Un skewness de +2.5 est fortement positif. Cela signifie que la distribution des salaires est \textbf{asymétrique à droite}.
Concrètement : la grande majorité des employés a un salaire regroupé (autour de la médiane), mais il existe une "longue queue" de quelques individus avec des salaires très élevés. Ces valeurs extrêmes "tirent" la moyenne vers la droite (Moyenne > Médiane).
\end{correctionbox}

\begin{correctionbox}[Correction Exercice 3 : Interprétation (Skewness)]
Si une distribution est parfaitement symétrique (comme la loi normale ou la loi uniforme), les écarts positifs et négatifs à la moyenne s'annulent parfaitement. Le skewness est \textbf{nul (0)}.
\end{correctionbox}

\begin{correctionbox}[Correction Exercice 4 : Interprétation (Kurtosis)]
1.  Un excess kurtosis > 0 est dit \textbf{leptokurtique}.
2.  Cela signifie que la distribution a des "queues plus épaisses" que la loi normale. La probabilité d'événements extrêmes (très grands gains ou très grandes pertes, comme un "krach") est \textbf{plus élevée} que ce qu'un modèle normal ne prédirait.
\end{correctionbox}

\begin{correctionbox}[Correction Exercice 5 : Interprétation (Kurtosis)]
1.  Un excess kurtosis < 0 est dit \textbf{platykurtique}.
2.  Visuellement, par rapport à une loi normale, la distribution est plus "aplatie" ou "carrée". Elle a un pic central moins prononcé et des queues plus "légères" (les événements extrêmes sont moins probables). La distribution uniforme est un exemple classique.
\end{correctionbox}

\begin{correctionbox}[Correction Exercice 6 : Comparaison (Skewness)]
La distribution exponentielle (ex: temps d'attente) a un pic à 0 et une longue queue vers les grandes valeurs. Elle est \textbf{asymétrique à droite}, et son skewness est \textbf{positif} (Skew = 2).
\end{correctionbox}

\begin{correctionbox}[Correction Exercice 7 : Comparaison (Kurtosis)]
La distribution uniforme est "plate" et n'a pas de pic central ni de queues épaisses. Elle est \textbf{platykurtique}. Son "excess kurtosis" est \textbf{négatif} (Kurtosis Excessif = -1.2).
\end{correctionbox}

\begin{correctionbox}[Correction Exercice 8 : Propriété de Symétrie]
Soit $Y = X - \mu$. La symétrie implique que $Y$ et $-Y$ (qui est $\mu - X$) ont la même distribution.
$E[(X-\mu)^3] = E[Y^3]$.
Puisque $Y$ et $-Y$ ont la même distribution, ils ont les mêmes moments : $E[Y^k] = E[(-Y)^k]$ pour tout $k$.
Pour $k=3$ :
$E[Y^3] = E[(-Y)^3] = E[(-1)^3 Y^3] = E[-Y^3] = -E[Y^3]$.
La seule valeur qui est égale à son opposé est 0.
$E[Y^3] = -E[Y^3] \implies 2 E[Y^3] = 0 \implies E[Y^3] = 0$.
Donc, $E[(X-\mu)^3] = 0$.
\end{correctionbox}

\begin{correctionbox}[Correction Exercice 9 : Vérification de Symétrie (PDF)]
On teste la condition $f(x) = f(2\mu - x)$. Le candidat pour $\mu$ est 3.
$f(2\mu - x) = f(2(3) - x) = f(6 - x)$.
$f(6 - x) = \frac{1}{2}e^{-|(6-x)-3|} = \frac{1}{2}e^{-|3-x|}$.
Puisque $|3-x| = |-(x-3)| = |x-3|$, on a $f(6-x) = \frac{1}{2}e^{-|x-3|} = f(x)$.
Oui, la distribution est symétrique autour de $\mu=3$.
\end{correctionbox}

\begin{correctionbox}[Correction Exercice 10 : Symétrie et Moments]
Non. Le Skewness nul est une condition nécessaire pour la symétrie, but pas suffisante.
On peut construire des distributions (assez exotiques) qui sont non symétriques, mais où les asymétries se compensent d'une manière qui annule le 3ème moment, résultant en un skewness de 0. Cependant, dans la plupart des cas pratiques, Skewness = 0 est un très bon indicateur de symétrie.
\end{correctionbox}

\begin{correctionbox}[Correction Exercice 11 : Définitions (Échantillon vs Pop.)]
$\sigma^2$ (variance de population) est un \textbf{paramètre} théorique. C'est la "vraie" variance de l'ensemble de la population (souvent inconnue).
$s^2$ (variance d'échantillon) est une \textbf{statistique}. C'est une \textit{estimation} de $\sigma^2$ calculée à partir d'un sous-ensemble de données (l'échantillon).
\end{correctionbox}

\begin{correctionbox}[Correction Exercice 12 : Correction de Bessel ($n-1$)]
Si $n=1$ (un seul échantillon $X_1$), la moyenne d'échantillon est $\bar{X} = X_1$.
La somme des carrés des écarts est $\sum (X_i - \bar{X})^2 = (X_1 - X_1)^2 = 0$.
Si on divisait par $n=1$, on obtiendrait $s^2 = 0/1 = 0$. Cela estimerait faussement que la population n'a pas de variance, ce qui est absurde.
Diviser par $n-1$ (donc $1-1=0$) donne $0/0$, une forme indéfinie, ce qui nous dit correctement qu'on ne peut pas estimer une dispersion à partir d'un seul point.
\end{correctionbox}

\begin{correctionbox}[Correction Exercice 13 : Calcul (Moments d'Échantillon)]
Échantillon : $\{ 2, 3, 10 \}$. $n=3$.
1.  $\bar{X} = \frac{1}{3} (2 + 3 + 10) = \frac{15}{3} = 5$.
2.  $s^2 = \frac{1}{n-1} \sum (X_i - \bar{X})^2 = \frac{1}{2} \left[ (2-5)^2 + (3-5)^2 + (10-5)^2 \right]$
    $s^2 = \frac{1}{2} \left[ (-3)^2 + (-2)^2 + (5)^2 \right]$
    $s^2 = \frac{1}{2} [ 9 + 4 + 25 ] = \frac{38}{2} = 19$.
\end{correctionbox}

\begin{correctionbox}[Correction Exercice 14 : MGF (Propriété de base)]
$M_X(t) = E[e^{tX}]$.
En $t=0$, $M_X(0) = E[e^{0 \cdot X}] = E[e^0] = E[1] = 1$.
Toute MGF valide doit valoir 1 en $t=0$.
\end{correctionbox}

\begin{correctionbox}[Correction Exercice 15 : MGF (Série de Taylor)]
On identifie les coefficients de la série $M_X(t) = 1 + \frac{E[X]}{1!}t + \frac{E[X^2]}{2!}t^2 + \dots$
1.  Le coefficient de $t$ est $E[X]$. On lit $5t$. Donc $E[X] = 5$.
2.  Le coefficient de $t^2$ est $E[X^2]/2!$. On lit $14t^2$.
    $E[X^2] / 2 = 14 \implies E[X^2] = 28$.
3.  $\text{Var}(X) = E[X^2] - (E[X])^2 = 28 - (5)^2 = 28 - 25 = 3$.
\end{correctionbox}

\begin{correctionbox}[Correction Exercice 16 : MGF (Calcul de moments)]
$M_X(t) = \lambda (\lambda - t)^{-1}$.
1.  $M_X'(t) = \lambda \cdot (-1) (\lambda - t)^{-2} \cdot (-1) = \lambda (\lambda - t)^{-2}$.
    $E[X] = M_X'(0) = \lambda (\lambda - 0)^{-2} = \lambda (\lambda^{-2}) = 1/\lambda$.
2.  $M_X''(t) = \lambda \cdot (-2) (\lambda - t)^{-3} \cdot (-1) = 2\lambda (\lambda - t)^{-3}$.
    $E[X^2] = M_X''(0) = 2\lambda (\lambda - 0)^{-3} = 2\lambda (\lambda^{-3}) = 2/\lambda^2$.
\end{correctionbox}

\begin{correctionbox}[Correction Exercice 17 : MGF (Variance de l'Exponentielle)]
$\text{Var}(X) = E[X^2] - (E[X])^2$
$\text{Var}(X) = (2/\lambda^2) - (1/\lambda)^2 = 2/\lambda^2 - 1/\lambda^2 = 1/\lambda^2$.
\end{correctionbox}

\begin{correctionbox}[Correction Exercice 18 : MGF (Loi Normale)]
$M_X(t) = \exp(\mu t + \frac{1}{2}\sigma^2 t^2)$.
On utilise la règle de la chaîne : $(\exp(u))' = \exp(u) \cdot u'$.
$M_X'(t) = \exp(\mu t + \frac{1}{2}\sigma^2 t^2) \cdot \frac{d}{dt}(\mu t + \frac{1}{2}\sigma^2 t^2)$
$M_X'(t) = \exp(\mu t + \frac{1}{2}\sigma^2 t^2) \cdot (\mu + \sigma^2 t)$.
$E[X] = M_X'(0) = \exp(0) \cdot (\mu + 0) = 1 \cdot \mu = \mu$.
\end{correctionbox}

\begin{correctionbox}[Correction Exercice 19 : MGF (Loi Normale)]
On dérive $M_X'(t)$ (règle du produit $u \cdot v$) :
$u = \exp(\mu t + \frac{1}{2}\sigma^2 t^2) \implies u' = u \cdot (\mu + \sigma^2 t)$
$v = (\mu + \sigma^2 t) \implies v' = \sigma^2$
$M_X''(t) = u'v + uv'$
$M_X''(t) = [\exp(\dots)(\mu + \sigma^2 t)] \cdot (\mu + \sigma^2 t) + [\exp(\dots)] \cdot (\sigma^2)$
On évalue en $t=0$ :
$E[X^2] = M_X''(0) = [\exp(0)(\mu)] \cdot (\mu) + [\exp(0)] \cdot (\sigma^2)$
$E[X^2] = (1 \cdot \mu) \cdot \mu + 1 \cdot \sigma^2 = \mu^2 + \sigma^2$.
\end{correctionbox}

\begin{correctionbox}[Correction Exercice 20 : MGF (Loi Normale)]
$\text{Var}(X) = E[X^2] - (E[X])^2$
$\text{Var}(X) = (\mu^2 + \sigma^2) - (\mu)^2 = \sigma^2$.
La MGF confirme bien que le paramètre $\sigma^2$ est la variance.
\end{correctionbox}

\begin{correctionbox}[Correction Exercice 21 : Propriété des Sommes]
Si $X$ et $Y$ sont indépendantes, la MGF de la somme est le \textbf{produit} des MGF :
$M_{X+Y}(t) = M_X(t) \cdot M_Y(t)$.
\end{correctionbox}

\begin{correctionbox}[Correction Exercice 22 : Application (Somme de Poissons)]
1.  $M_S(t) = M_X(t) \cdot M_Y(t) = e^{\lambda_1(e^t - 1)} \cdot e^{\lambda_2(e^t - 1)}$
    $M_S(t) = e^{\lambda_1(e^t - 1) + \lambda_2(e^t - 1)} = e^{(\lambda_1 + \lambda_2)(e^t - 1)}$.
2.  On reconnaît la MGF d'une loi de Poisson. Le paramètre est ce qui multiplie $(e^t - 1)$.
    Donc $S$ suit une loi de Poisson de paramètre $(\lambda_1 + \lambda_2)$.
    $S \sim \text{Poisson}(\lambda_1 + \lambda_2)$.
\end{correctionbox}

\begin{correctionbox}[Correction Exercice 23 : Application (Somme de Binomiales)]
$M_S(t) = M_X(t) \cdot M_Y(t) = (p e^t + (1-p))^n \cdot (p e^t + (1-p))^m$
$M_S(t) = (p e^t + (1-p))^{n+m}$.
On reconnaît la MGF d'une loi Binomiale avec $n+m$ essais et une probabilité de succès $p$.
$S \sim \text{Bin}(n+m, p)$.
\end{correctionbox}

\begin{correctionbox}[Correction Exercice 24 : Application (Transformation Linéaire)]
$M_Y(t) = E[e^{tY}] = E[e^{t(aX + b)}] = E[e^{taX + tb}]$
$M_Y(t) = E[e^{taX} \cdot e^{tb}]$.
Puisque $e^{tb}$ est une constante :
$M_Y(t) = e^{tb} E[e^{(ta)X}]$.
Par définition, $E[e^{(ta)X}]$ est la MGF de $X$ évaluée au point $(ta)$.
$M_Y(t) = e^{tb} M_X(ta)$.
\end{correctionbox}

\begin{correctionbox}[Correction Exercice 25 : Application (Moyenne d'Échantillon)]
Soit $S = \sum X_i$. $\bar{X} = S / n = (1/n)S$.
1.  D'abord, la MGF de la somme $S$ (Exercice 21) :
    $M_S(t) = M_{X_1}(t) \cdot \dots \cdot M_{X_n}(t) = [M_X(t)]^n$ (car i.i.d.)
2.  Ensuite, on utilise la transformation linéaire $\bar{X} = aS + b$ avec $a=1/n$ et $b=0$ (Exercice 24) :
    $M_{\bar{X}}(t) = e^{b t} M_S(a t) = e^0 M_S(t/n)$
    $M_{\bar{X}}(t) = M_S(t/n)$.
3.  On combine les deux :
    $M_{\bar{X}}(t) = [M_X(t/n)]^n$.
\end{correctionbox}
